% Separate research note. Failed openness inference only; not a new status row.
\subsection*{A jumping-boundary test for the proposed open-fullness shortcut}

This note concerns the injectivity condition in a finite-cycle
realization theorem. It does not give a counterexample to finite
generation, and no representation-count conclusion is asserted.

Let $k$ be algebraically closed of characteristic different from two.
Choose a smooth projective curve $C$ of genus two and a Weierstrass
point $P$, so that $\omega_C\simeq\mathcal O_C(2P)$.
Put $N=\mathcal O_C(-3P)$ and let $Y=\operatorname{Tot}(N)$,
with zero section again denoted by $C$ and projection $r:Y\to C$.
The argument takes place on the thickenings of this zero section.
Write $C_j=V(\mathcal I_C^j)$.

The reduced curve $C$ is a weak reduction cycle:
$\mathcal O_C(-C)=\mathcal O_C(3P)$ is generically generated and
has vanishing first cohomology. It is \emph{not} a strong reduction
cycle, since its intrinsic dualizing inverse has negative degree.
To meet the strong-cycle hypotheses, fix any $a\geq5$ and put
$Z=aC$. Indeed,
\[
 \mathcal O_Z(-Z)=r^*\mathcal O_C(3aP)|_Z,
 \qquad
 \omega_Z^{-1}=r^*\mathcal O_C((3a-5)P)|_Z.
\]
Here $\omega_Y=r^*\mathcal O_C(5P)$ and adjunction gives the
second equality. Both displayed lines are globally generated:
their degrees on $C$ are at least $2g(C)=4$, and their generating
sections pull back along $r$. Moreover
\[
 H^1(Z,\mathcal O_Z(-Z))
   =\bigoplus_{j=0}^{a-1}H^1(C,\mathcal O_C(3(a+j)P))=0,
\]
using the split algebra
$\mathcal O_Z=\bigoplus_{j=0}^{a-1}N^{-j}$ of the line-bundle
neighbourhood. Thus $Z$ has the strong reduction-cycle properties
used in the realization packet.

Let $T\subset C$ be an affine neighbourhood of $P$, with variable
point $Q$. On $C\times T$ put
\[
 L=\mathcal O_C(4P),\qquad
 M_Q=\mathcal O_C(P+Q),\qquad V_Q=L\oplus M_Q.
\]
The notation $M_Q$ denotes the line bundle
$\mathcal O_{C\times T}(P\times T+\Delta)$, where $\Delta$
is the relative diagonal. Every $V_Q$ is generically generated.
For every $j\geq1$,
\[
                 H^1(C,V_Q\otimes N^{-j})=0.              \tag{B.NO.1}
\]
This follows componentwise from the degrees $4+3j$ and $2+3j$.

\paragraph{An actual family on the first ribbon.}
The first ribbon is split as an algebra,
$\mathcal O_{C_2}=\mathcal O_C\oplus N^{-1}$ with square-zero
second summand. Bundle lifts of $V$ to $C_2\times T$, with its
reduction specified, have the usual Cech deformation classes in
$H^1(C\times T,\mathcal End(V)\otimes N^{-1})$.
Use the off-diagonal summand from $L$ to $M_Q$:
\[
 \mathcal Hom(L,M_Q)\otimes N^{-1}=\mathcal O_C(Q).
                                                               \tag{B.NO.2}
\]
For every $Q$, both $h^0\mathcal O_C(Q)$ and
$h^1\mathcal O_C(Q)$ equal one. Consequently its relative first
cohomology is a line bundle on $T$ and commutes with base change.
After shrinking $T$, choose a section $\epsilon$ that is nonzero
at $P$ (and hence everywhere on a further neighbourhood).
Because $T$ is affine, this is an actual relative Cech class.
It defines a bundle $F_2$ on $C_2\times T$: lift the transition
matrices of $V$ and multiply them by $1+\epsilon_{ij}$ in the
off-diagonal square-zero ideal. The cocycle equation is precisely
the additive Cech cocycle equation, so this is a genuine bundle
family, not only fibrewise extension data.

Its boundary map in
\[
 0\longrightarrow V_Q\longrightarrow F_{2,Q}(C)
     \longrightarrow V_Q\otimes N\longrightarrow0
\]
has domain
\[
 H^0(C,V_Q\otimes N)
   =H^0(C,\mathcal O_C(P))\oplus
      H^0(C,\mathcal O_C(Q-2P))=k.                        \tag{B.NO.3}
\]
Its only possible nonzero component is multiplication by the
section of $\mathcal O_C(P)$, applied to the class (B.NO.2):
\[
 H^1(C,\mathcal O_C(Q))
       \longrightarrow H^1(C,\mathcal O_C(P+Q)).          \tag{B.NO.4}
\]
At $Q=P$, this map is nonzero. By Serre duality its dual is the
nonzero map $H^0(C,\mathcal O_C)\to H^0(C,\mathcal O_C(P))$,
and both spaces have dimension one. Thus the boundary (B.NO.3)
is injective at $P$. For $Q\ne P$ it is not injective, because
\[
 H^1(C,\mathcal O_C(P+Q))
       =H^0(C,\mathcal O_C(P-Q))^\vee=0.                 \tag{B.NO.5}
\]
A degree-zero line with a section is trivial, and two distinct
points on a positive-genus curve are not linearly equivalent.

\paragraph{The same failure on the strong cycle $Z=aC$.}
Extend $F_2$ successively to a bundle $F_{2a}$ on $C_{2a}\times T$.
At each step the obstruction lies in second coherent cohomology
of a curve over the affine base $T$, hence is zero. No uniqueness
of this extension is needed. Put $V_{a,Q}=F_{2a,Q}|_{C_a}$.
The positive conormal layers in (B.NO.1) show that
\[
 H^1(C_a,V_{a,Q})\simeq H^1(C,V_Q),\qquad
 H^0(C_a,V_{a,Q})\twoheadrightarrow H^0(C,V_Q).            \tag{B.NO.6}
\]
In particular $V_{a,Q}$ is generically generated on the full
nonreduced cycle. Furthermore $F_{2a,Q}(-aC)$ is globally
generated on $C_{2a}$ for every $Q$: its reduced restriction
has summands of degrees $4+3a$ and $2+3a$, so it is globally
generated; all subsequent conormal layers have vanishing $H^1$,
so the generating sections lift, and nilpotent Nakayama finishes.
Thus this family is compatible with the open generating-tuple
charts used for a fixed negative twist.

For every $j\geq2$, one has $H^0(C,V_Q\otimes N^j)=0$.
The pole-layer filtration therefore identifies
\[
 H^0(C_a,V_{a,Q}(aC))\simeq H^0(C,V_Q\otimes N)=k.
                                                               \tag{B.NO.7}
\]
Under (B.NO.6)--(B.NO.7), the boundary of
\[
 0\longrightarrow V_{a,Q}\longrightarrow F_{2a,Q}(aC)
    \longrightarrow V_{a,Q}(aC)\longrightarrow0
\]
is exactly the first-ribbon boundary (B.NO.4). To see the
compatibility directly, include $V_Q(C)$ as the highest
nilpotent layer of $V_{a,Q}(aC)$, take its preimage in
$F_{2a,Q}(aC)$, and push out the kernel along
$V_{a,Q}\to V_Q$. The resulting extension is
$0\to V_Q\to F_{2,Q}(C)\to V_Q(C)\to0$.
Its boundary is consequently injective at $Q=P$ and fails
to be injective at every nearby $Q\ne P$.

\paragraph{Exact scope of the obstruction.}
Even on a genuine strong reduction cycle, and even when a fixed
negative twist of every object is globally generated, the
field-valued boundary-injectivity condition need not be open.
It holds at the special point in this family and fails at its
generic point. The first cohomology in (B.NO.6) jumps from one
to zero, so this family does not satisfy the coflat hypotheses
of the qualified Artinian realization theorem. Those hypotheses
cannot be removed merely by declaring fullness to be an open
derived-vanishing condition. This does not rule out an alternative
finite-cycle Frobenius realization argument with additional
equations or witnesses.
